Nuclear enrichment involves separating uranium isotopes, primarily by increasing the concentration of uranium-235 (U-235) in uranium hexafluoride (UF6) from its natural level of about 0.7 percent to levels suitable for nuclear fuel, typically between 3 and 5 percent for most commercial reactors. In the United States, nuclear power generates approximately 18 to 19 percent of the nation's electricity, serving as a reliable baseload source that produces more carbon-free power than any other energy type, with a capacity factor exceeding 92 percent, making it one of the cleanest and most energy-dense options available.
The cost of uranium enrichment is primarily determined by two factors: the energy efficiency of the process, measured in separative work units (SWU), and the throughput capacity for processing bulk material. SWU quantifies the effort required to separate isotopes as a function of feed, product, and tails concentrations and is independent of the specific technology used. Over the past 50 years, enrichment methods have evolved from gaseous diffusion to gas centrifuges, with laser-based techniques emerging as a third-generation approach promising greater efficiency. Among these, laser enrichment, such as the Separation of Isotopes by Laser Excitation (SILEX) process, demonstrates the highest SWU efficiency.
Laser enrichment employs an infrared laser tuned to a specific frequency to selectively excite the molecular vibrations of UF6 containing U-235, allowing its separation from U-238 using a collector mechanism. This method is estimated to be significantly more energy-efficient than gaseous diffusion, with costs potentially as low as $30 per SWU compared to higher figures for older technologies, representing a substantial improvement over diffusion's energy demands.
The price of enriched uranium increases with higher U-235 concentrations, as the separative work required scales logarithmically with the diminishing mass processed at each stage. For instance, enriching to 5 percent requires far less effort than to 20 percent, and prices per SWU reflect this non-linear progression. Uranium enriched to 20 percent or above is classified as high-assay low-enriched uranium (HALEU), which is not weapons-grade—typically defined as 90 percent or higher U-235—and is unavailable for commercial markets due to proliferation concerns. In early 2026, spot prices for uranium oxide concentrate ranged from $80 to $94 per pound, while SWU costs averaged about $98 in recent years.
Global Laser Enrichment (GLE), a joint venture between Australia's Silex Systems (51 percent) and Canada's Cameco (49 percent), holds exclusive rights to the SILEX laser enrichment technology and has been advancing its commercialization, achieving Technology Readiness Level 6 in 2025 and securing $28 million in U.S. Department of Energy funding for further development. In the past, they have had enormous challenges scaling their technology, but they are planning a new facility in Padukah, Kentucky.

The market for 3 percent enriched uranium, used in conventional light-water reactors, is driven by stable demand from approximately 94 operational reactors across 54 plants in the United States. These plants feature experienced procurement teams that treat enriched uranium as a dependable commodity, with refueling cycles occurring every 18 to 24 months. Much of this material has been sourced from abroad, particularly Russia, but the Prohibiting Russian Uranium Imports Act, signed in May 2024 and effective from August 2024 through 2040, restricts such imports, with waivers possible until 2028 to prevent supply disruptions.
High-assay low-enriched uranium (HALEU), enriched to between 5 and 19.75 percent U-235, commands different pricing due to its targeting of startups and advanced reactor projects, which often have less experienced buyers and constrained capital. Demand for HALEU is expanding as advanced reactors require it for efficient operation, but supply remains limited, prompting U.S. investments, such as $2.7 billion in enrichment contracts, to bolster domestic production. HALEU production costs are estimated at around $23,725 to $25,725 per kilogram, reflecting its specialized nature.
New market entrants in the enrichment market are aiming for the cheapest enrichment process. In a recent interview, Scott Nolan from General Matter said the main lesson he learned from his experience at SpaceX was how to keep lowering the cost of a kilogram to low Earth orbit. The analog for the nuclear industry is lowering the cost per swu of HALEU.
Laser-based enrichment technologies pose proliferation risks because their efficiency enables smaller, more concealable facilities that could accelerate the production of weapons-grade uranium with fewer detectable signatures. Consequently, they are subject to stringent regulations, including export controls and safeguards imposed by bodies such as the International Atomic Energy Agency, to mitigate the potential for misuse by state or non-state actors. The federal government still enriches nuclear material for its weapons program. There have been many claims that the work done at sites like PNNL in Washington has had difficulty bridging the demands of the commercial and government sectors.
